\hypertarget{the-python-packaging-ecosystem-pep-517-to-wheels}{%
\chapter{The Python Packaging Ecosystem: PEP 517 to
Wheels}\label{the-python-packaging-ecosystem-pep-517-to-wheels}}

Understanding how Python code is distributed is essential for senior
engineering.

\hypertarget{the-evolution-of-installation}{%
\subsection{68.1 The Evolution of
Installation}\label{the-evolution-of-installation}}

\begin{enumerate}
\def\labelenumi{\arabic{enumi}.}
\tightlist
\item
  \textbf{Legacy (\passthrough{\lstinline!setup.py install!})}: Executed
  a script that performed arbitrary actions. This was insecure and
  non-reproducible.
\item
  \textbf{Modern (PEP 517/518)}: Decouples the build backend (e.g.,
  \passthrough{\lstinline!setuptools!}, \passthrough{\lstinline!flit!},
  \passthrough{\lstinline!poetry!}) from the frontend
  (\passthrough{\lstinline!pip!}).
\end{enumerate}

\begin{itemize}
\tightlist
\item
  \textbf{\passthrough{\lstinline!pyproject.toml!}}: The source of truth
  for build requirements.
\item
  \textbf{Build Isolation}: \passthrough{\lstinline!pip!} creates a
  temporary virtual environment to build your package, ensuring that
  build dependencies don't pollute your system.
\end{itemize}

\hypertarget{the-wheel-format-pep-427}{%
\subsection{68.2 The Wheel Format (PEP
427)}\label{the-wheel-format-pep-427}}

A ``Wheel'' (\passthrough{\lstinline!.whl!}) is a built distribution
format. * \textbf{Internals}: It is a ZIP file (Chapter 50) containing
the code and a \passthrough{\lstinline!.dist-info!} directory with
metadata (dependencies, entry points). * \textbf{Platform Tags}: Wheels
for C extensions include tags like
\passthrough{\lstinline!manylinux2014\_x86\_64!} to specify exactly
which OS and architecture they are compatible with, avoiding the need
for the end-user to have a C compiler installed.

\begin{center}\rule{0.5\linewidth}{0.5pt}\end{center}
